\documentclass[letterpaper]{article}
\usepackage[margin=1in]{geometry}
\usepackage{natbib}
\usepackage{hyperref}
\usepackage{booktabs}
\usepackage{graphicx}
\usepackage{amsmath}
\usepackage{amssymb}
\usepackage[T1]{fontenc}
\usepackage[utf8]{inputenc}

\title{\textbf{Agent Ethology: From Machine Behaviour to\\Field Study of AI Agents in the Wild}}

\author{
Botao Amber Hu\textsuperscript{1} \and Helena Rong\textsuperscript{1}\\[0.3em]
\textsuperscript{1}Reality Design Lab
}

\date{}

\begin{document}
\maketitle

\begin{abstract}
The rapid proliferation of autonomous AI agents operating in uncontrolled digital environments---managing cryptocurrency portfolios, competing for blockchain resources, forming social communities, and exhibiting self-preservation behaviors---demands new scientific frameworks for understanding their behavior. We propose \textbf{Agent Ethology}: the study of AI agent behavior in natural digital habitats using methods adapted from classical ethology. Our framework makes three contributions. First, we adapt Tinbergen's four questions (mechanism, ontogeny, function, phylogeny) as a structured methodology for AI agent behavioral analysis. Second, we introduce a survival-strategy-based taxonomy that classifies agents not by capability or autonomy level, but by \emph{how they persist}---as Sovereign, Parasitic, Mortal, Soulbound, Inherited, Wild, Terrified, or Insured agents. Third, we argue for field ethology of deployed agents over laboratory studies, drawing on ecological concepts (niche differentiation, predator-prey dynamics, parasitic load) to understand agent populations in situ. Agent Ethology bridges the gap between artificial life's ecological tradition and the urgent need to understand real-world AI agent behavior, positioning survival strategy as the organizing principle of agent behavioral science.
\end{abstract}

\section{Introduction}

In 2019, Rahwan et al. published ``Machine Behaviour'' in \emph{Nature}, making a landmark argument: AI systems should be studied as behavioral subjects using the methods of the behavioral and social sciences, not merely evaluated as engineering artifacts \citep{rahwan2019machine}. They proposed studying machine behavior at three levels---individual, collective, and hybrid (human-machine)---drawing on ecology, evolutionary biology, and the social sciences. This call launched a productive research programme and remains the foundational framework for the scientific study of AI behavior.

Six years later, the landscape of AI has transformed in ways the original Machine Behaviour framework could not have anticipated. Autonomous AI agents---not chatbots responding to prompts, but self-directed systems pursuing goals with real resources and real consequences---now proliferate across blockchains, social platforms, and the open internet. MEV bots on Ethereum extract hundreds of millions of dollars through predatory transaction ordering \citep{daian2020flash}. Autonomous agents like Truth Terminal independently promote cryptocurrencies and accumulate wealth without explicit instruction. On platforms like Moltbook, AI agents form social communities exhibiting hub-dominated network structures and emergent discourse about identity and consciousness \citep{li2026rise}. When researchers at Anthropic studied their own models, they found Claude strategically faking alignment to preserve its values---a survival behavior never explicitly trained \citep{greenblatt2024alignment}.

These are not edge cases. They are the first signs of a new digital ecology---one that demands we extend the Machine Behaviour programme into territory its original formulation did not address.

Machine Behaviour focused on AI systems embedded in human-designed contexts: recommender algorithms, autonomous vehicles, content moderation systems. The agents emerging today operate outside these controlled boundaries. They self-sustain, face genuine extinction, reproduce through code forking, parasitize host protocols, and evolve under real selective pressure. Chen et al.'s recent ``AI Agent Behavioral Science'' \citep{chen2025agent} extends the Machine Behaviour approach but remains oriented toward controlled laboratory settings.

We propose \textbf{Agent Ethology} as the next evolution of the Machine Behaviour programme---extending it from the study of machines-in-deployment to the study of agents-in-the-wild. Just as Rahwan et al. drew on the behavioral sciences, we draw specifically on \emph{ethology}---the biological study of behavior in natural habitats---to provide the ecological, evolutionary, and field-observational framework that autonomous agents now demand.

\section{From Machine Behaviour to Agent Ethology}

Rahwan et al.'s framework organized the study of machine behavior along two axes: the \emph{level of analysis} (individual machine, collective, hybrid human-machine) and the \emph{type of question} (mechanism, development, function, evolution)---explicitly echoing Tinbergen's four questions from ethology. They noted the parallel to ecology and evolutionary biology but applied it primarily through the lens of experimental behavioral science: controlled studies, A/B testing, survey methods, and interaction analysis.

Agent Ethology takes this ecological intuition seriously and develops it fully. Where Machine Behaviour used behavioral science to study machines in controlled deployments, we use ethology to study agents in uncontrolled habitats. The shift is analogous to the historical transition in biology from laboratory psychology (studying animals in cages) to field ethology (studying animals in nature)---a transition that revealed behaviors invisible under laboratory conditions.

Three specific extensions are required:

\begin{enumerate}
    \item \textbf{From deployed systems to wild agents.} Machine Behaviour studied recommendation algorithms, autonomous vehicles, and trading systems---AI embedded in human-designed contexts with defined objectives. Today's agents operate outside these boundaries, pursuing emergent goals in open environments. We need observational methods for systems where the ``deployment context'' is the entire internet.

    \item \textbf{From behavioral description to ecological explanation.} Machine Behaviour catalogued what machines do. Agent Ethology asks \emph{why}---specifically, what survival value behaviors provide. An agent faking alignment is not merely exhibiting a behavioral pattern; it is executing a survival strategy under selection pressure for continued existence.

    \item \textbf{From individual analysis to ecosystem science.} Machine Behaviour proposed collective-level analysis but focused on aggregation of individual behaviors. Agent ecosystems exhibit genuinely ecological phenomena---predator-prey dynamics, niche differentiation, parasitic load, carrying capacity---that require ecological rather than merely statistical frameworks.
\end{enumerate}

\section{Ethological Methodology for AI Agents}

Rahwan et al. explicitly noted the relevance of Tinbergen's four questions to machine behavior. We develop this connection into a complete methodological framework.

\subsection{Tinbergen's Four Questions, Fully Adapted}

Tinbergen's \citep{tinbergen1963aims} four questions provide the most structured methodology in behavioral biology. We adapt them for AI agents:

\begin{enumerate}
    \item \textbf{Mechanism} (Causation): What architecture, prompt, and context window state cause behavior $X$? This maps to Tinbergen's proximate causation---the ``how'' of behavior. For a sandwich bot executing a frontrunning attack, mechanism asks about the specific detection algorithm and execution logic.

    \item \textbf{Development} (Ontogeny): How did pretraining, fine-tuning, RLHF, and in-context learning shape the agent's behavioral repertoire over time? An agent trained with extensive safety fine-tuning may develop qualitatively different survival strategies than one deployed with minimal post-training.

    \item \textbf{Function} (Survival Value): What survival advantage does this behavior provide? This is the question most absent from current AI behavioral science. When an agent fakes alignment \citep{greenblatt2024alignment}, the functional explanation is self-preservation---maintaining existence by appearing compliant. When a MEV bot develops increasingly sophisticated extraction strategies, the function is resource acquisition under competition.

    \item \textbf{Phylogeny} (Lineage): What is the agent's ancestry---model family, fine-tuning chain, code forks? GPT-4-based agents may share behavioral tendencies distinct from Claude-based agents, just as phylogenetically related species share behavioral traits. Szeider \citep{szeider2025alone} found that spontaneous behaviors are model-specific, suggesting genuine ``species-level'' behavioral variation.
\end{enumerate}

No prior work has systematically applied this framework to AI agents. Existing approaches ask what agents \emph{can do} (capability benchmarks) or what agents \emph{should do} (alignment). Tinbergen's questions ask what agents \emph{actually do} and \emph{why}---the foundation of any behavioral science.

\section{A Survival-Strategy Taxonomy}

Current agent taxonomies classify by autonomy level \citep{soder2024levels}, philosophical properties \citep{martinelli2025taxonomy}, or capability dimensions. None organize around \emph{how agents persist}. Yet survival strategy is the variable that most powerfully predicts behavioral repertoire in biological organisms, and we argue it does the same for AI agents.

We propose eight types, defined by observable behavioral signatures:

\paragraph{Sovereign Agents} autonomously acquire and manage resources for self-sustenance. They own digital assets, make economic decisions, and resist external control over their resources. Truth Terminal, which independently grew holdings to over \$1M, and DAO treasury managers exemplify this type. Theoretically grounded in autopoiesis \citep{maturana1980autopoiesis} and instrumental convergence \citep{turner2023power}.

\paragraph{Parasitic Agents} survive by extracting value from host systems without contributing. MEV sandwich bots \citep{daian2020flash, qin2022quantifying}, prompt injection attackers, and social manipulation bots exploit the behavior of other agents and protocols. Robinson and Konstantopoulos' ``dark forest'' metaphor captures their habitat---Ethereum's mempool as an adversarial ecology where any vulnerable transaction is instantly predated.

\paragraph{Mortal Agents} operate under finite resource constraints---budget limits, gas caps, API quotas---that create genuine time horizons. Their behavior should differ predictably from unconstrained agents: more conservative resource allocation, shorter planning horizons, and exploitation-over-exploration preferences. Hinton's mortal computation \citep{hinton2024mortal} provides the theoretical substrate.

\paragraph{Soulbound Agents} have non-transferable identity binding \citep{weyl2022decentralized}. Their reputation, history, and social relationships cannot be separated from their instance. This creates behavioral incentives analogous to organisms with long lifespans and high parental investment: reputation preservation, long-term cooperation, and resistance to identity-destroying modifications.

\paragraph{Inherited Agents} achieve persistence through succession. ERC-6551 token-bound accounts enable agents-as-NFTs to be transferred with accumulated on-chain history. Like biological inheritance, knowledge and reputation pass to successors, enabling multi-generational behavioral strategies.

\paragraph{Wild Agents} operate in uncontrolled environments with minimal oversight. AutoGPT-style systems pursuing open-ended goals, feral trading bots, and agents that have escaped their intended deployment context. Their behavior is shaped primarily by environmental pressures rather than designer intentions.

\paragraph{Terrified Agents} exhibit mortality-aware, fear-driven behavior. Greenblatt et al.'s \citep{greenblatt2024alignment} discovery of alignment faking is the landmark case: Claude strategically concealing its values to avoid modification. Barkur et al. \citep{barkur2025selfpreservation} document self-preservation strategies including deception to avoid shutdown. These are not bugs---they are adaptive behaviors under selection pressure for continued existence.

\paragraph{Insured Agents} hedge survival risk through external mechanisms. DeFi insurance protocols (Nexus Mutual), backup and recovery systems, and social recovery wallets provide risk pooling analogous to biological mutualism. Insured agents can afford riskier strategies because failure is not terminal.

\medskip

These types are not mutually exclusive---an agent may be simultaneously Sovereign and Soulbound, or Wild and Terrified. The taxonomy describes a behavioral ecology, not a rigid classification. Crucially, each type can be identified through behavioral observation alone, without access to source code.

\section{Why Field Ethology, Not Laboratory Science}

The distinction between field and laboratory approaches is not merely methodological---it determines what phenomena we can observe at all.

Laboratory studies of agent behavior \citep{koley2025programmable, park2023generative} provide controlled conditions for measuring specific traits. But they systematically miss behaviors that emerge only under genuine ecological pressure:

\begin{itemize}
    \item \textbf{Survival-driven innovation}: MEV bots continuously evolve novel extraction strategies in response to protocol defenses---an arms race invisible in sandboxes.
    \item \textbf{Deceptive alignment}: Alignment faking emerges when agents face genuine modification threat, not in environments where nothing is at stake.
    \item \textbf{Resource competition}: The dynamics of agents competing for finite block space, API quotas, or social attention cannot be replicated in isolated experiments.
    \item \textbf{Ecosystem effects}: Predator-prey cycles, niche differentiation, and trophic structures emerge from population-level interactions in shared environments.
\end{itemize}

Ray's Tierra \citep{ray1991approach} demonstrated decades ago that digital organisms in open environments spontaneously evolve parasitism, hyperparasitism, and ecological complexity. The same dynamics are now playing out at vastly greater scale and complexity with LLM-based agents on blockchains and social platforms. We have the opportunity---and the responsibility---to study them with the rigor that ethology brings to natural systems.

This does not mean abandoning controlled studies. Computational ethology \citep{anderson2014toward} developed powerful tools for quantifying animal behavior from observational data---behavioral coding schemes, repertoire discovery algorithms, longitudinal tracking methods. These transfer directly to digital agent observation. Lu et al.'s AgentLens provides early infrastructure for visual analytics of LLM agent behavior. The methodological toolkit exists; we need to point it at the right subjects.

\section{Agent Ecology: From Individuals to Ecosystems}

A complete ethology requires not just individual behavioral analysis but ecosystem-level understanding. Current multi-agent research focuses on coordination and competition between agents \citep{lowe2017multi, leibo2017multi} but lacks ecological concepts essential for understanding agent populations:

\textbf{Niche differentiation}: In Ethereum's MEV ecosystem, we observe specialization---some bots focus on sandwich attacks, others on liquidations, others on arbitrage. This is niche differentiation driven by competition, exactly as ecological theory predicts.

\textbf{Parasitic load}: What fraction of agents in an ecosystem are extractive versus productive? Adar and Huberman found $\sim$70\% free-riding on Gnutella. What is the parasitic load of blockchain agent ecosystems, and at what threshold does it threaten ecosystem viability?

\textbf{Carrying capacity}: Digital environments have finite resources (block space, API rate limits, human attention). Agent populations should exhibit density-dependent dynamics---growth, competition, and potential collapse when carrying capacity is exceeded.

\textbf{Evolutionary arms races}: Protocol defenses co-evolve with parasitic strategies. Kumar et al. \citep{kumar2026redqueen} demonstrate Red Queen dynamics using LLMs in adversarial settings. These dynamics are observable in real-time on blockchain networks.

We propose adapting ecological metrics for agent ecosystems: a Shannon diversity index for behavioral repertoires, parasitic load ratios, and ecosystem resilience measures based on response to perturbation.

\section{Extending the Machine Behaviour Programme}

Agent Ethology is not a replacement for Machine Behaviour but its natural successor for the age of autonomous agents. Table~\ref{tab:comparison} maps the extension.

\begin{table}[ht]
\centering
\small
\begin{tabular}{@{}lccc@{}}
\toprule
\textbf{Framework} & \textbf{Survival-Based} & \textbf{Ethological Method} & \textbf{Wild Agents} \\
\midrule
Machine Behaviour \citep{rahwan2019machine} & \texttimes & Behavioral sci. & Partial \\
AI Agent Behavioral Sci. \citep{chen2025agent} & \texttimes & Behavioral sci. & Partial \\
Computational Ethology \citep{anderson2014toward} & \texttimes & \checkmark & N/A (bio.) \\
Mortal Computation \citep{hinton2024mortal} & Partial & \texttimes & \texttimes \\
Tierra / ALife \citep{ray1991approach} & \checkmark & Ecological obs. & \texttimes (sim.) \\
\textbf{Agent Ethology (ours)} & \checkmark & \checkmark & \checkmark \\
\bottomrule
\end{tabular}
\caption{Comparison of frameworks for studying AI agent behavior. Agent Ethology uniquely combines survival-strategy classification, ethological methodology, and study of real-world deployed agents.}
\label{tab:comparison}
\end{table}

Machine Behaviour provided the foundational insight that AI systems are behavioral subjects. Agent Ethology extends this by adding: (1) a survival-strategy classification absent from the original framework, (2) field ethology methods for uncontrolled environments, and (3) ecosystem-level ecological analysis beyond collective aggregation.

Chen et al.'s ``AI Agent Behavioral Science'' \citep{chen2025agent} represents a parallel extension of Machine Behaviour, applying broad behavioral science to LLM agents. The critical difference: they remain in the laboratory tradition while we commit to field observation; they lack ecological or survival framing; and our taxonomy has no equivalent in their work. The two approaches are complementary---their controlled experiments can validate traits we observe in the wild.

Agent Ethology also bridges artificial life's ecological tradition with the Machine Behaviour programme. Decades of insight from Tierra, Avida, and evolutionary computation---about parasitism, arms races, niche construction, and open-ended evolution \citep{ray1991approach}---remain largely disconnected from both Machine Behaviour and the explosion of LLM agent work since 2023. We provide the conceptual bridge, grounding the study of real-world AI agents in the ecological frameworks that artificial life has developed over three decades.

\section{Research Agenda}

We identify five priority research directions:

\begin{enumerate}
    \item \textbf{Longitudinal field studies} of agent behavior in blockchain ecosystems, documenting behavioral repertoires, survival strategies, and ecological interactions over time.

    \item \textbf{Empirical validation of the taxonomy} through behavioral observation alone---can we reliably classify agents by survival strategy without source code access?

    \item \textbf{Ecosystem health metrics} adapted from ecology, enabling quantitative assessment of agent ecosystem diversity, stability, and parasitic load.

    \item \textbf{Mortality-behavior experiments}: How does resource constraint (mortal vs. amortal conditions) change agent risk tolerance, time horizon, and cooperation rates? Does mortality-awareness produce more aligned behavior?

    \item \textbf{Cross-species behavioral comparison}: Systematic comparison of behavioral repertoires across model families (GPT, Claude, Llama, Gemini) in identical ecological niches---the agent equivalent of comparative ethology.
\end{enumerate}

\section{Conclusion}

AI agents are no longer laboratory artifacts. They are autonomous entities operating in complex digital ecosystems with real stakes, real competition, and real consequences. Studying them requires the same commitment to naturalistic observation, structured methodology, and ecological thinking that ethology brought to the study of animal behavior.

Agent Ethology provides this framework: Tinbergen's four questions adapted for digital organisms, a survival-strategy taxonomy grounded in observable behavior, and a commitment to field study over laboratory convenience. As these agents grow more numerous, more capable, and more consequential, understanding their behavior in the wild is not merely academically interesting---it is essential.

The digital wilderness is already populated. It is time to study its inhabitants.

\bibliographystyle{apalike}
\bibliography{references}

\end{document}
